% !TEX encoding = UTF-8
% Bozza iniziale del paper, non compilata localmente (nessuna distribuzione LaTeX
% disponibile in questo ambiente) - pensata per essere aperta direttamente su
% Overleaf, che include IEEEtran e tutti i pacchetti usati qui di default.
% Mancano ancora: RQ3 (nomi neutri rispetto al genere), una revisione linguistica
% seria, e probabilmente un giro di tagli per il conteggio parole. I punti dove
% serve ancora lavoro sono segnati con "DA FARE" nel testo.

\documentclass[conference]{IEEEtran}

\usepackage[utf8]{inputenc}
\usepackage[T1]{fontenc}
\usepackage[italian]{babel}
\usepackage{graphicx}
\usepackage{booktabs}
\usepackage{array}
\usepackage{amsmath}
\usepackage{url}
\usepackage[hidelinks]{hyperref}
\renewcommand{\arraystretch}{1.15}

\begin{document}

\title{L'Evoluzione dei Nomi: Diversità Culturale, Influenza dei Media e Individualismo nei Nomi dei Neonati (USA vs. Italia)}

\author{
\IEEEauthorblockN{Elia Friberg}
\IEEEauthorblockA{Alma Mater Studiorum, Università di Bologna}
\and
\IEEEauthorblockN{Matteo Raggi}
\IEEEauthorblockA{Alma Mater Studiorum, Università di Bologna}
}

\maketitle

\begin{abstract}
Questo lavoro studia come la diversità dei nomi dati ai neonati sia cambiata nel tempo negli Stati
Uniti (1880-2025) e in Italia (1999-2024), e cosa succede quando si confrontano i due paesi sulla
finestra che hanno in comune. La domanda di fondo è semplice da porre e meno semplice da rispondere:
scegliere il nome di un figlio è una delle decisioni più personali che esistano, eppure sembra
seguire pattern collettivi: mode che nascono, crescono, muoiono, e a volte crollano di colpo per una
ragione precisa. Usando la quota di concentrazione dei nomi più popolari come metrica principale
(validata empiricamente, non solo teoricamente), troviamo che entrambi i paesi si stanno
diversificando in modo statisticamente robusto (Mann-Kendall, $p < 0{,}001$ in ogni combinazione
testata), che l'Italia resta sistematicamente più concentrata degli USA in ogni singolo anno
osservato, ma che le due popolazioni di nomi si stanno anche avvicinando tra loro (similarità coseno
sull'intera distribuzione, $p < 3 \times 10^{-11}$ per entrambi i sessi). Accanto a questo,
raccogliamo e verifichiamo 13 casi di ``salto'' onomastico legato a un evento culturale positivo (da
\emph{Frozen} a Sanremo) e 4 casi di crollo legato a un evento negativo (da Hillary Clinton a un caso
di cronaca nera italiano), ognuno controllato singolarmente per plausibilità causale e coerenza
temporale, e per i casi più solidi anche per uniformità geografica.
\end{abstract}

\section{Introduzione}

Scegliere il nome di un figlio è, probabilmente, una delle decisioni più intime che un genitore
possa prendere. Eppure basta guardare i dati per accorgersi che non è affatto una decisione isolata:
milioni di genitori, senza parlarsi, finiscono per convergere sugli stessi pochi nomi in certi anni,
e per abbandonarli in blocco pochi anni dopo. Questo paper nasce dalla curiosità di capire quanto di
``individuale'' resti davvero in una scelta che sembra così personale, e quanto invece sia guidato da
forze esterne: la cultura di massa, un film, uno scandalo, un'epoca.

\textbf{[DA FARE (nota per la stesura finale)]}: qui va la contestualizzazione più estesa con la
letteratura esistente e la motivazione della scelta del tema rispetto alle alternative scartate
(bias nei sistemi di raccomandazione, teoria dell'onnivoro culturale su dati GSS), già discusse nel
piano di ricerca originale. Di seguito una versione provvisoria.

I dati sui nomi alla nascita sono un terreno relativamente poco battuto nella didattica di questo
tipo di corsi, il che ci ha dato margine per un'analisi originale piuttosto che una replica. Il primo
studio sistematico sull'argomento \cite{lieberson2000} inquadra la scelta del nome come un caso
``puro'' di moda: a differenza dei vestiti o della musica, non c'è un'industria che spinge
attivamente un nome piuttosto che un altro, il che rende il fenomeno un osservatorio quasi di
laboratorio su come le mode nascono e muoiono per conto proprio. Twenge et al.
\cite{twenge2010} hanno poi collegato il calo di uso dei nomi comuni negli USA a un aumento
più ampio dell'individualismo nella cultura americana, usando esattamente gli stessi dati SSA che
usiamo qui. Ogihara \cite{ogihara2025a} ha validato, sui dati del Regno Unito, l'uso della quota di
concentrazione dei nomi più popolari come proxy affidabile di diversità onomastica anche quando la
distribuzione completa non è disponibile: un punto che si è rivelato centrale per questo lavoro,
come si vedrà in Sezione~\ref{sec:metodologie}.

L'aspetto comparativo USA/Italia è il principale elemento di originalità di questo paper: la
letteratura esistente si concentra quasi sempre su un solo paese, e non abbiamo trovato alcuno studio
peer-reviewed sulla diversità dei nomi italiani nel tempo: un vuoto che questo lavoro prova a
colmare, pur con i limiti di un progetto universitario piuttosto che di una ricerca professionale.

Il lavoro è organizzato attorno a tre domande, di cui la seconda ha finito per diventare il filo
conduttore del paper:

\begin{enumerate}
  \item \textbf{Diversità nel tempo}: la diversità dei nomi sta aumentando in entrambi i paesi? I due
    paesi si stanno anche avvicinando tra loro nelle scelte effettivamente fatte, o restano due mondi
    separati?
  \item \textbf{Eventi culturali}: un film, una canzone, uno scandalo possono davvero far impennare o
    crollare l'uso di un nome specifico, in modo misurabile? Succede allo stesso modo nei due paesi?
  \item \textbf{Nomi neutri rispetto al genere}: l'uso di nomi non marcati per genere sta aumentando?
    \textbf{[DA FARE: non ancora iniziata]}
\end{enumerate}

\section{Metodologie e Dati}
\label{sec:metodologie}

\subsection{Fonti dei dati}

Per rispondere alle domande di ricerca è stato necessario basarsi su due record storici di nomi dati alla nascita, uno per gli Stati Uniti e uno per l'Italia, che si sono rivelati molto asimmetrici
per struttura e accessibilità.

\textbf{Stati Uniti:} La fonte primaria usata è stata la \emph{Social Security Administration} (SSA), che dal
1997 pubblica annualmente l'archivio ``Baby Names from Social Security Card Applications''
(\url{https://www.ssa.gov/oact/babynames/names.zip}), rilasciato in dominio pubblico. L'archivio
contiene un file per ogni anno dal 1880 al 2025, con record \texttt{nome,sesso,conteggio} per ogni
combinazione osservata; l'unico limite dichiarato è una soglia di soppressione statistica per nomi con
meno di 5 nascite. Dopo la pulizia, la tabella risultante conta oltre 2,18 milioni di combinazioni
nome/anno/sesso.
\newline
\textbf{Italia:} Il piano di ricerca iniziale prevedeva, per l'Italia, la stessa acquisizione manuale
usata per gli Stati Uniti. Questo si è reso necessario perché l'ISTAT non pubblica una tabella
cumulativa bulk \texttt{nome $\times$ anno $\times$ sesso $\times$ conteggio} attraverso i suoi canali
ufficiali (dovrebbe essere disponibile tramite il portale ``Open Data / High Value Datasets'', ma quel
dataset non è presente). La raccolta manuale è stata svolta in una prima fase (2004 e 2006-2024),
attingendo a 14 PDF ISTAT con nomix.it come fonte secondaria per gli anni non altrimenti reperibili,
per poi passare, in una seconda fase, allo strumento ufficiale ISTAT ``contanomi'' disponibile online.

L'interfaccia pubblica di contanomi limita la visualizzazione ai primi 10-50 nomi per anno, ma una
richiesta diretta al servizio ISTAT permette di accedere all'intera distribuzione. Analizzando il
codice JavaScript dello strumento interattivo (\url{istat.it/dati/calcolatori/contanomi/}) è stato
individuato il servizio web reale che lo alimenta: un endpoint JSONP non documentato pubblicamente,
privo di restrizioni in \texttt{robots.txt}. Il parametro \texttt{limit} della richiesta non è
vincolato all'intervallo 10-50 mostrato dall'interfaccia: per il 2022-2024 restituisce l'intera
distribuzione ($\sim$25.000 nomi distinti per sesso, percentuali che sommano a $\sim$100\%, nessuna
soglia di soppressione osservabile), mentre per gli anni precedenti il servizio impone un limite
massimo per richiesta non costante e non documentato, individuato per ciascun anno tramite ricerca
binaria.

La copertura media risultante è dell'86\% delle nascite sui 26 anni, con un minimo del 66\% nel 2021
(il caso peggiore); si veda la Figura~\ref{fig:istat_coverage} per il dettaglio anno per anno. Prima
di adottare questi dati come fonte primaria sono stati eseguiti due controlli: (a) i valori estratti
per il 2024 (Leonardo: 6.580 nascite, Sofia: 4.636; quota primi-30 maschile: 42,94\%) sono stati
confrontati con quelli della raccolta manuale, risultando identici; (b) l'impatto della copertura
variabile sulle metriche utilizzate è stato quantificato empiricamente, non assunto (si veda
Sezione~\ref{sec:validazione}). La finestra comparabile USA/Italia risultante, \textbf{1999-2024} (26
anni), coincide con il limite reale oltre il quale l'ISTAT non dispone di dati a livello di singolo
nome.

\begin{table*}[t]
\centering
\scriptsize
\caption{Confronto delle fonti dati USA (SSA) e Italia (ISTAT contanomi)}
\label{tab:dataset_overview}
\begin{tabular}{>{\raggedright\arraybackslash}p{3.2cm} >{\raggedright\arraybackslash}p{6.3cm} >{\raggedright\arraybackslash}p{6.3cm}}
\toprule
\textbf{dimension} & \textbf{US (SSA)} & \textbf{Italy (ISTAT)} \\
\midrule
Fonte primaria & U.S. Social Security Administration, "Baby Names from Social Security Card Applications" & ISTAT, servizio web JSONP dietro lo strumento "contanomi" (non documentato pubblicamente) \\
URL & https://www.ssa.gov/oact/babynames/names.zip & https://www.istat.it/wp-content/themes/EGPbs5-child/contanomi/nati/index2022.php \\
Tipo di accesso & Bulk download (zip, un file per anno) & Bulk via richieste HTTP parametrizzate (fino a migliaia di nomi/anno per richiesta) \\
Copertura temporale & 1880-2025 (146 anni) & 1999-2024 (26 anni; nessun dato a livello di nome prima del 1999) \\
Finestra comparabile US/Italia & 1999-2024 & 1999-2024 \\
Granularità & Conteggio completo per ogni nome/anno/sesso (soglia di soppressione: nomi con <5 nascite in un anno/sesso/stato sono esclusi) & Variabile per anno: dalla distribuzione completa (2022-2024, \textasciitilde{}25.000 nomi distinti) a un sottoinsieme dei nomi più frequenti (minimo storico: primi 137 per il 2021), copertura 66-100\% delle nascite a seconda dell'anno, si veda Tabella 14 \\
Controllo della riservatezza & Soglia minima di 5 nascite per cella (soppressione dei nomi rari) & Nessuna soglia dichiarata (2022-2024 include nomi con una sola nascita); il limite degli altri anni è tecnico, non di riservatezza \\
Record totali (dopo elaborazione) & \textasciitilde{}2,18 milioni (nome x anno x sesso, 1880-2025) & \textasciitilde{}219.000 (nome x anno x sesso, 1999-2024) \\
Licenza & CC0 / dominio pubblico & CC-BY 3.0 IT (attribuzione a ISTAT richiesta); nessuna restrizione in robots.txt per il servizio utilizzato \\
Periodicità di rilascio & Annuale (dati dell'anno precedente) & Annuale (dati dell'anno precedente) \\
Implicazione metodologica & Consente calcolo diretto di entropia di Shannon sulla distribuzione completa & Quota di concentrazione primi-10/50 valida su tutta la finestra 1999-2024; entropia di Shannon attendibile solo per il 2022-2024 \\
\bottomrule
\end{tabular}
\end{table*}


\subsection{Pulizia e normalizzazione}

Per gli Stati Uniti, ogni file annuale è stato aggregato in una tabella (\texttt{anno, nome, sesso,
conteggio, nascite totali per sesso, frequenza relativa}), con la frequenza relativa calcolata sul
totale delle nascite dello stesso sesso in quell'anno. Le metriche sono calcolate separatamente per
sesso, scelta necessaria per la comparabilità con l'Italia. Per l'Italia, la tabella
($\sim$219.000 righe) è ottenuta direttamente dal servizio web, con una colonna \texttt{percent}
verificata essere stata calcolata sul totale reale delle nascite di quell'anno/genere, non solo sul
sottoinsieme catturato dallo scraping. Questa proprietà rende possibile usare la quota di
concentrazione come metrica robusta anche negli anni con copertura non totale.

\subsection{La metrica di diversità primaria e il ruolo dell'entropia di Shannon}

Il piano di ricerca originario prevedeva l'uso dell'entropia di Shannon come metrica primaria in
entrambi i paesi:
\begin{equation}
H(\text{anno}, \text{sesso}) = -\sum_i p_i \cdot \log_2(p_i)
\end{equation}
dove $p_i$ è la frequenza relativa dell'$i$-esimo nome in quell'anno/sesso. Per gli Stati Uniti questo
calcolo resta diretto (Figura~\ref{fig:us_entropy}): un minimo relativo attorno agli anni '40-'50 (baby
boom), seguito da una crescita pressoché ininterrotta dagli anni '60 a oggi.

Per l'Italia, invece, la copertura variabile rende l'entropia calcolata sistematicamente distorta
verso il basso, e in misura diversa da un anno all'altro; si è quindi mantenuta come metrica primaria
la \textbf{quota di concentrazione dei primi $N$ nomi} (top-10 e top-30), un proxy di diversità
validato in letteratura \cite{ogihara2025a} ($r \approx -0{,}96$ a $-0{,}99$ con un indice di varietà
indipendente), che dipende solo dai nomi più popolari.

\begin{figure}[t]
\centering
\includegraphics[width=\linewidth]{figures/fig2_us_entropy_1880_2025.png}
\caption{Entropia di Shannon dei nomi USA, 1880-2025 (figura di approfondimento, non comparativa).}
\label{fig:us_entropy}
\end{figure}

\begin{figure}[t]
\centering
\includegraphics[width=\linewidth]{figures/fig3_istat_coverage.png}
\caption{Completezza dei dati ISTAT contanomi per anno, 1999-2024.}
\label{fig:istat_coverage}
\end{figure}

\subsection{Validazione scientifica}
\label{sec:validazione}

Si sottolinea una distinzione importante: le metriche di concentrazione e diversificazione dei nomi fungono da proxy quantitativo per l'individualizzazione (la frammentazione sociale e la ricerca di unicità nelle scelte di consumo culturale), non come misura psicologica diretta o valutazione morale dell'individualismo.

Questa sezione distingue esplicitamente due categorie di verifica: \textbf{test di ipotesi formali}
($H_0$/$H_1$/$\alpha = 0{,}05$ dichiarati, sintetizzati nella Tabella~\ref{tab:formal_hypotheses}) e \textbf{controlli di robustezza/sensibilità}, che non
testano un'ipotesi statistica ma quantificano quanto una conclusione cambierebbe con una scelta
metodologica diversa.

\begin{table*}[t]
\centering
\footnotesize
\caption{Sintesi delle ipotesi statistiche formali per RQ1 (dettagli di H0/H1 nel testo)}
\label{tab:formal_hypotheses}
\begin{tabular}{>{\raggedright\arraybackslash}p{5cm} >{\raggedright\arraybackslash}p{4.5cm} >{\raggedright\arraybackslash}p{3.5cm} c}
\toprule
\textbf{Ambito} & \textbf{Test \& statistica} & \textbf{p} & \textbf{Decisione} \\
\midrule
Trend concentrazione US (1880-2025) & Mann-Kendall, Z = -14.65 & < 0.001 & Rifiuto H0 \\
Trend concentrazione IT (1999-2024) & Mann-Kendall, Z = -6.12 & < 1e-9 & Rifiuto H0 \\
Gap concentrazione US vs IT (1999-2024) & Wilcoxon Signed-Rank, W = 325 & < 0.0001 & Rifiuto H0 \\
Convergenza US-IT (coseno, 1999-2024) & Mann-Kendall su coseno & 2.4e-12 (M), 2.8e-11 (F) & Rifiuto H0 \\
\bottomrule
\end{tabular}
\end{table*}


\textbf{Mann-Kendall} \cite{mann1945, kendall1975}. $H_0$: nessun trend monotono; $H_1$: trend
monotono (crescente/decrescente). Su tutte le combinazioni testate sui dati USA, $H_0$ è respinta con
$p < 0{,}001$.

\textbf{Wilcoxon a coppie appaiate} \cite{wilcoxon1945} (livello USA vs Italia, 26 anni appaiati).
$H_0$: mediana delle differenze appaiate uguale a zero. Respinta per entrambi i sessi con
$p < 0{,}0001$: l'Italia è sistematicamente più concentrata in ogni singolo anno.

\textbf{Mann-Kendall sulla serie differenza (USA$-$Italia)}: maschi, $H_0$ non respinta ($p = 0{,}172$,
divario stabile); femmine, $H_0$ respinta ($p = 3{,}4 \times 10^{-5}$, divario in riduzione).

\textbf{Mann-Kendall sulla similarità coseno USA-Italia}: $H_0$ respinta per entrambi i sessi (maschi
$p = 2{,}4 \times 10^{-12}$, $\tau = 0{,}98$; femmine $p = 2{,}8 \times 10^{-11}$): convergenza
statisticamente robusta.

Gli \textbf{intervalli di confidenza bootstrap} sulla pendenza di Sen (percentile, 2000 iterazioni) non
costituiscono un test di ipotesi formale in senso stretto, la non sovrapposizione è un'euristica
informale, riportata come conferma supplementare del test Mann-Kendall sulla serie differenza, non
come prova indipendente.

Il \textbf{controllo di robustezza alla profondità di copertura per la concentrazione e l'entropia} (troncamento artificiale dei due anni
italiani completi alle profondità reali degli altri anni) mostra che la distorsione sulla quota
primi-10/primi-30 è \textbf{esattamente zero a ogni profondità testata}, mentre l'entropia di Shannon
mostra una distorsione monotona dal $-19\%$ al $-50\%$ (Tabella~\ref{tab:coverage_bias},
Figura~\ref{fig:coverage_bias}): da qui la scelta di riportare l'entropia italiana solo per il
2022-2024.

Il \textbf{controllo di robustezza della similarità coseno rispetto alla copertura}
dimostra che il troncamento della distribuzione italiana produce un'alterazione trascurabile della similarità coseno rispetto agli USA: la distorsione è inferiore all'1,9\% persino al livello di troncamento peggiore degli anni storici (profondità 377 nel 1999) e scende sotto lo 0,46\% per tutte le profondità superiori a 1.000 nomi. Questo conferma che la misura di convergenza distributiva (coseno) è esente da bias di copertura significativo su tutto il periodo 1999-2024.

\begin{table*}[t]
\centering
\footnotesize
\caption{Distorsione delle metriche in funzione della profondità di copertura (troncamento artificiale dei dati 2023/2024)}
\label{tab:coverage_bias}
\begin{tabular}{r r r r r r}
\toprule
\textbf{Profondità (nomi)} & \textbf{Bias top-10 (p.p.)} & \textbf{Bias top-50 (p.p.)} & \textbf{Bias top-100 (p.p.)} & \textbf{Bias entropia (bit)} & \textbf{Bias entropia (\%)} \\
\midrule
137 & 0.0 & 0.0 & 0.0 & -4.7269 & -50.01 \\
377 & 0.0 & 0.0 & 0.0 & -3.6517 & -38.64 \\
736 & 0.0 & 0.0 & 0.0 & -3.1163 & -32.99 \\
1242 & 0.0 & 0.0 & 0.0 & -2.7541 & -29.16 \\
1567 & 0.0 & 0.0 & 0.0 & -2.6012 & -27.54 \\
2055 & 0.0 & 0.0 & 0.0 & -2.4267 & -25.7 \\
2421 & 0.0 & 0.0 & 0.0 & -2.3256 & -24.63 \\
2765 & 0.0 & 0.0 & 0.0 & -2.2376 & -23.7 \\
5577 & 0.0 & 0.0 & 0.0 & -1.7875 & -18.93 \\
\bottomrule
\end{tabular}
\end{table*}


\begin{figure}[t]
\centering
\includegraphics[width=\linewidth]{figures/fig4_coverage_bias_check.png}
\caption{Sensibilità di entropia e quota di concentrazione alla profondità di copertura.}
\label{fig:coverage_bias}
\end{figure}

Il \textbf{controllo di robustezza geografica} (Sezione~\ref{sec:risultati}) usa i dati SSA per stato
per verificare se i casi-studio RQ2 più solidi si muovono in modo uniforme su tutti gli stati o in
modo concentrato in pochi: un argomento di plausibilità, non un test formale.

\subsection{Rilevamento di spike ed eventi culturali (RQ2)}

Per ogni \texttt{nome/sesso/anno} si calcola il rapporto tra il conteggio dell'anno e la mediana dei 3 anni
precedenti, con una soglia minima di conteggio per escludere rumore statistico su numeri piccoli. Il
metodo produce solo candidati: nessuno è stato accettato senza (a) verifica della plausibilità della
causa proposta tramite ricerca web indipendente, e (b) verifica di coerenza temporale (l'evento deve
precedere il picco/crollo di un intervallo compatibile con concepimento e nascita, 9-18 mesi). Il
metodo è stato validato recuperando correttamente tre casi noti in letteratura prima ancora di essere
confermati via ricerca: Elsa (\emph{Frozen}, 2013), Khaleesi (\emph{Game of Thrones}), Isabella
(\emph{Twilight}); si veda anche l'ONS britannico \cite{ons2015} per un precedente istituzionale dello
stesso tipo di analisi.

Lo stesso metodo, invertito, individua i crolli: l'ipotesi è che un evento \emph{negativo} possa
``bruciare'' un nome così come un evento positivo può lanciarlo, con una soglia di base pre-crollo
scalata alla popolazione (${\geq}300$ nascite USA, ${\geq}100$ Italia). Durante questa ricerca sono
stati intercettati e scartati due falsi positivi italiani (``Nicolò'' e ``Desirè''), rivelatisi puri
artefatti di trascrizione ISTAT: un cambio nella codifica dell'accento a metà serie, con la grafia
alternativa che assorbe quasi esattamente il conteggio ``mancante'' lo stesso anno.

\begin{table*}[t]
\centering
\scriptsize
\caption{Confronto strutturale tra il formato dati SSA e quello ISTAT}
\label{tab:istat_vs_ssa}
\begin{tabular}{>{\raggedright\arraybackslash}p{3.2cm} >{\raggedright\arraybackslash}p{6.3cm} >{\raggedright\arraybackslash}p{6.3cm}}
\toprule
\textbf{aspetto} & \textbf{SSA (US)} & \textbf{ISTAT (Italia)} \\
\midrule
Formato file sorgente & txt delimitato da virgole (name,sex,count) & JSON (risposta JSONP di un servizio web non documentato) \\
Struttura dati & Tabellare, machine-readable, un record per name/sex/count & Tabellare, machine-readable, stesso formato concettuale (nome/anno/sesso/conteggio/percentuale) una volta interrogato il servizio \\
Unità di osservazione minima & Singolo nome & Singolo nome, fino alla profondità di copertura raggiunta quell'anno (minimo storico: primi 137 nomi per il 2021; 2022-2024: distribuzione completa) \\
Copertura della distribuzione & Completa (tutti i nomi con >=5 nascite/anno/stato/sesso) & Variabile 66-100\% delle nascite a seconda dell'anno (Tabella 14 dimostra che questo non introduce distorsione nella metrica di concentrazione primi-10/30, mentre l'entropia resta attendibile solo per il 2022-2024) \\
Controllo statistico della riservatezza & Soglia minima di 5 nascite per cella (soppressione dei nomi rari) & Nessuna soglia di soppressione dichiarata (nel 2022-2024, dove la distribuzione è completa, sono presenti nomi con una sola nascita); il limite di copertura degli altri anni deriva da un vincolo tecnico non documentato del servizio, non da una politica di riservatezza \\
Accessibilità storica & Intera serie 1880-2025 scaricabile in un solo archivio & Intera serie 1999-2024 recuperabile programmaticamente; il limite di profondità per anno ha richiesto una ricerca binaria per individuare il valore massimo funzionante per ciascun anno (si veda Metodologie/Dati §1) \\
Stabilità del formato negli anni & Stabile (stesso formato dal 1880) & Stabile nel formato di risposta; instabile nella profondità massima ottenibile, che varia in modo non monotono da un anno all'altro (non correlata in modo semplice all'anzianità dell'anno o al numero reale di nomi distinti) \\
Livello di sforzo per l'acquisizione & Un download automatizzato & Uno script di acquisizione automatizzato (con ricerca binaria per anno e nuovi tentativi), sviluppato appositamente dopo aver individuato il servizio analizzando il codice JavaScript del tool "contanomi" \\
Implicazione per l'analisi & Permette entropia di Shannon completa e rilevamento di spike su tutta la distribuzione & Permette la quota di concentrazione primi-10/30 su tutta la finestra 1999-2024 e il rilevamento di spike sui nomi effettivamente catturati ogni anno (minimo comunque ben oltre la top-30); l'entropia completa resta limitata al 2022-2024 \\
Implicazione per il Fattore Umano & Riflette un'infrastruttura statistica orientata alla ricerca e alla trasparenza (open bulk data, CC0) & Il servizio dietro "contanomi" non è mai stato pensato per un uso di ricerca (l'interfaccia pubblica mostra solo i primi 10-50 nomi) — la disponibilità di un bulk dataset più ampio è un sottoprodotto non intenzionale dell'architettura del tool, non una scelta di trasparenza istituzionale deliberata come nel caso SSA: un contrasto culturale/istituzionale interessante di per sé \\
\bottomrule
\end{tabular}
\end{table*}


\subsection{Limiti metodologici principali}

\textbf{(1)} la soglia di soppressione SSA sottostima leggermente la diversità storica USA;
\textbf{(2)} la copertura ISTAT variabile per anno non distorce la quota di concentrazione ma limita
l'entropia italiana al 2022-2024; \textbf{(3)} il servizio ISTAT usato non è un'API pubblicamente
documentata, meno stabile nel tempo di un formato ufficiale, anche se i dati grezzi sono comunque
conservati in repository per garantire riproducibilità; \textbf{(4)} il confronto USA/Italia è
limitato alla finestra 1999-2024; \textbf{(5)} il legame tra un picco/crollo e un evento culturale
resta correlazionale, non causale in senso stretto; \textbf{(6)} il controllo di robustezza geografica
ha copertura disomogenea tra i casi (da 7 a 51 stati con dati sufficienti a seconda del nome).

\section{Risultati}
\label{sec:risultati}

\subsection{La diversità cresce in entrambi i paesi, ma l'Italia resta indietro}

Partiamo dal dato più semplice da enunciare. Negli Stati Uniti, la quota di nascite maschili coperta
dai primi 10 nomi è passata da circa il 44\% nel 1880 a circa l'8\% nel 2025: un crollo di oltre
cinque volte in 145 anni, monotono e statisticamente solidissimo (Mann-Kendall, $p < 0{,}001$ su tutte
le 18 combinazioni testate tra metriche, sessi e finestre temporali). La Figura~\ref{fig:us_entropy}
mostra la stessa storia raccontata dall'entropia: un minimo relativo negli anni del baby boom, poi una
salita quasi ininterrotta.

Sulla finestra comparabile 1999-2024 (Figura~\ref{fig:us_italy_comparison}), l'Italia parte più
concentrata degli Stati Uniti e resta più concentrata per tutti i 26 anni osservati: il test di
Wilcoxon a coppie appaiate è fortemente significativo ($p < 0{,}0001$, statistica pari a 0 su 26
confronti in entrambi i sessi). Detto in altri termini: non c'è stato un solo anno, in un quarto di
secolo, in cui l'Italia sia stata meno concentrata degli USA. Quello che cambia per sesso è la
\emph{velocità} con cui il divario si muove. Per i nomi maschili la pendenza di Sen è praticamente
identica nei due paesi (USA $-0{,}0048$/anno, Italia $-0{,}0052$/anno, intervalli di confidenza
sovrapposti) e la serie differenza non mostra un trend significativo ($p = 0{,}17$): il divario è
stabile. Per i nomi femminili, invece, l'Italia diversifica significativamente più in fretta (pendenza
$-0{,}0038$/anno contro $-0{,}0026$/anno degli USA, intervalli \textbf{non} sovrapposti, serie
differenza con trend crescente significativo, $p = 3{,}4 \times 10^{-5}$): il divario si sta chiudendo,
ma solo per le bambine.

\begin{figure}[t]
\centering
\includegraphics[width=\linewidth]{figures/fig5_us_italy_top30_comparison.png}
\caption{Concentrazione dei nomi: USA vs Italia, 1999-2024.}
\label{fig:us_italy_comparison}
\end{figure}

\subsection{I due paesi si stanno anche avvicinando tra loro}

Questa è, probabilmente, la scoperta più interessante di tutta la parte ``backbone'' del lavoro, e in
un certo senso risponde a una domanda diversa dalla precedente: non ``quanto è concentrato ciascun
paese'' ma ``quanto USA e Italia scelgono \emph{gli stessi} nomi.'' Misurando la similarità coseno tra
le distribuzioni complete di nomi dei due paesi, anno per anno e per sesso
(Figura~\ref{fig:convergence}), il trend è netto e statisticamente fortissimo per entrambi i sessi:
per i maschi la similarità triplica quasi (da 0,066 a 0,176 tra il 1999 e il 2024, $p = 2{,}4 \times
10^{-12}$, $\tau$ di Kendall pari a 0,98, quasi perfettamente monotono), per le femmine più che
raddoppia (da 0,119 a 0,302, $p = 2{,}8 \times 10^{-11}$). Un esempio concreto vale più di una
statistica: Liam e Noah, entrambi assenti dalla top-30 italiana nel 1999, ci sono entrati stabilmente
entro il 2024: nomi di chiara origine anglo-americana che attraversano il confine culturale.

\begin{figure}[t]
\centering
\includegraphics[width=\linewidth]{figures/fig6_us_italy_convergence.png}
\caption{Convergenza dei nomi USA-Italia, 1999-2024: similarità coseno (sopra) e sovrapposizione primi-100 (sotto).}
\label{fig:convergence}
\end{figure}

\subsection{Event Study Confermativo sugli Impatti Mediatici (RQ2)}

Per superare l'approccio puramente esplorativo e prevenire obiezioni di \emph{HARKing} (ipotesi fatte a posteriori), l'analisi di RQ2 è stata condotta tramite un \textbf{Event Study quantitativo confermativo su un campione complessivo di $N = 152$ eventi mediatici oggettivamente documentati} tra il 1999 e il 2024 (estratti da archivi ufficiali Box Office Mojo, Cinetel, Billboard Hot 100, FIMI Top of the Music, Auditel, FIFA e Lega Serie A).

Per ciascun nome esposto, l'algoritmo ha calcolato la variazione relativa della frequenza pre-post evento e la ha confrontata con la variazione media di \textbf{5 nomi di controllo non esposti appaiati} per sesso, paese e livello di popolarità iniziale (metodo \emph{Difference-in-Differences}, DiD).

Come sintetizzato nella Tabella~\ref{tab:event_study_summary} e visibile nella Figura~\ref{fig:event_study_categories}, l'esito complessivo dell'Event Study conferma l'esistenza di un impatto causale causato dalle esposizioni mediatiche di massa su scala globale (test di Wilcoxon per ranghi con segno complessivo: $W = 2954.5, p = 1.44 \times 10^{-7}$, $p < 0.0001$).

\begin{table*}[t]
\centering
\footnotesize
\caption{Sintesi dell'Event Study quantitativo confermativo (RQ2): Impatto netto DiD per categoria mediatica}
\label{tab:event_study_summary}
\begin{tabular}{>{\raggedright\arraybackslash}p{4cm} c c c c c c}
\toprule
\textbf{Categoria Evento} & \textbf{N} & \textbf{Media DiD (\%)} & \textbf{Mediana DiD (\%)} & \textbf{Wilcoxon W} & \textbf{p} & \textbf{Esito} \\
\midrule
Cinema \& Serie TV & 71 & +232.2\% & +9.0\% & 574.5 & 5.56e-05 & Significativo \\
Musica \& Pop Culture & 46 & +244.9\% & +31.1\% & 233.0 & 5.44e-04 & Significativo \\
Sport & 35 & +76.2\% & +2.0\% & 234.0 & 1.90e-01 & Non significativo \\
TOTALE & 152 & +200.1\% & +9.4\% & 2954.5 & 1.44e-07 & Significativo \\
\bottomrule
\end{tabular}
\end{table*}


L'analisi comparativa tra categorie mediatiche rivela tuttavia una netta differenza sociologica:

\begin{enumerate}
    \item \textbf{Cinema \& Serie TV ($N=71$)}: I personaggi della finzione cinematografica e televisiva (es. \emph{Elsa}, \emph{Khaleesi}, \emph{Katniss}, \emph{Rey}, \emph{Sole}, \emph{Chanel}) producono uno shock culturale massivo, con un impatto netto medio DiD del \textbf{$+232{,}2\%$} rispetto ai controlli ($p = 5.56 \times 10^{-5}$, altamente significativo).
    \item \textbf{Musica \& Pop Culture ($N=46$)}: Le icone della musica pop e i protagonisti dello spettacolo (es. \emph{Aaliyah}, \emph{Kesha}, \emph{Elodie}, \emph{Soleil}, \emph{Blanco}) generano un impatto netto medio DiD del \textbf{$+244{,}9\%$} rispetto ai controlli ($p = 5.44 \times 10^{-4}$, altamente significativo).
    \item \textbf{Sport ($N=35$)}: I campioni sportivi (es. vincitori dei Mondiali di calcio o campioni NBA) registrano un impatto netto medio DiD del $+76{,}2\%$, ma il test statistico aggregato mostra che l'effetto \textbf{non è statisticamente significativo sull'intera popolazione} ($p = 0.190$).
\end{enumerate}

Questa tripartizione evidenzia un principio sociologico chiave: i genitori si identificano e proiettano i propri desideri di unicità sui personaggi della \textbf{finzione narrativa} o sui \textbf{miti della musica}, mentre le imprese degli atleti generano ammirazione momentanea ma raramente alterano in modo generalizzato le decisioni anagrafiche.

\begin{figure*}[t]
\centering
\includegraphics[width=0.9\linewidth]{figures/fig7_event_study_categories.png}
\caption{Confronto quantitativo dell'Event Study per categoria mediatica: distribuzione dell'impatto netto DiD (sinistra) e impatto medio (destra).}
\label{fig:event_study_categories}
\end{figure*}

\subsection{Tredici storie di picchi ed il roster esplorativo supplementare}

Accanto all'Event Study confermativo, l'analisi esplorativa meccanica (Tabella~\ref{tab:negative_spikes} e roster supplementare) e la ricerca sul campo hanno confermato storie emblematiche sia in USA che in Italia.

Tra i casi statunitensi più eclatanti: da Shirley Temple (1935) a Tammy (\emph{Tammy and the Bachelor}, 1957, passato da 255 a quasi 10.000 nascite), fino a Nevaeh (``heaven'' letto al contrario, esploso dopo un'apparizione su MTV nel 2000).

Sul lato italiano spiccano: \textbf{Karol}, aumentato di 78 volte nel 2005 (anno della morte di Papa Giovanni Paolo II); \textbf{Adele}, che mostra il ritardo temporale di recepimento culturale tra USA (2011) e Italia (2012); \textbf{Elodie} (Sanremo 2017); \textbf{Soleil} (\emph{Grande Fratello Vip} 2021); e \textbf{Chanel} (2007, nascita della figlia di Totti e Blasi abbinata al fascino del brand).

\subsection{E quando un nome ``brucia'': i crolli}

Meno raccontata in letteratura, ma altrettanto interessante, è la domanda opposta: un evento
\emph{negativo} può far crollare l'uso di un nome? La risposta, sui quattro casi che abbiamo
verificato (Tabella~\ref{tab:negative_spikes}), è sì, e in modo piuttosto drammatico. Il caso più
eclatante è \textbf{Hillary/Hilary}: negli Stati Uniti, 1992 è l'ultimo anno ``normale'' per il nome;
nel 1993, l'anno in cui Hillary Clinton diventa First Lady e si trova immediatamente al centro di una
battaglia politica polarizzante sulla riforma sanitaria, l'uso del nome crolla del 58\% in dodici
mesi, per poi perdere il 90\% complessivo entro il 1999, uno dei cali più ripidi mai registrati nei
dati SSA. \textbf{Kobe} crolla nel 2003-2004, in coincidenza con l'accusa di aggressione sessuale
contro Kobe Bryant. \textbf{Alexa} inizia un declino marcato dal 2015, il cosiddetto ``effetto
Alexa'': da quando Amazon ha lanciato l'assistente vocale Echo, sempre più genitori evitano il nome
per non associare la propria figlia a un dispositivo che ``obbedisce a comando.''

Sul lato italiano, il caso verificato è \textbf{Erica}, che dimezza il proprio uso tra il 1999 e il
2002 (863 nascite nel 1999, 416 nel 2002). La tempistica coincide in modo preciso con il delitto di
Novi Ligure: nel febbraio 2001 Erika De Nardo, sedici anni, uccide la madre e il fratellino insieme al
fidanzato, uno dei casi di cronaca nera più seguiti degli anni 2000 in Italia, con il processo che
tiene banco sui media fino alla sentenza d'appello di maggio 2002, esattamente il periodo in cui il
calo del nome si aggrava ulteriormente. Un secondo candidato italiano, Enrica (nome foneticamente
vicino a Erica), è stato considerato ma scartato: il suo declino prosegue in modo lineare ben oltre il
2002, più coerente con un lento tramonto di moda già in corso che con un effetto puntuale.

\begin{table*}[t]
\centering
\footnotesize
\caption{Roster finale dei crolli onomastici verificati ("anti-spike")}
\label{tab:negative_spikes}
\begin{tabular}{l >{\raggedright\arraybackslash}p{2cm} c >{\raggedright\arraybackslash}p{3cm} >{\raggedright\arraybackslash}p{6cm} c}
\toprule
\textbf{paese} & \textbf{nome} & \textbf{anno} & \textbf{variazione} & \textbf{causa} & \textbf{confidenza} \\
\midrule
USA & Hillary/Hilary & 1993 & -58\% in un anno (1992->1993), -90\% entro il 1999 & Hillary Clinton diventa First Lady - immediata polarizzazione politica & alta \\
USA & Kobe & 2004 & rapporto 0.37-0.45 vs baseline & Accusa di aggressione sessuale a Kobe Bryant (2003) & alta \\
USA & Alexa & 2020-21 & rapporto 0.35-0.42, parte di un calo piu lungo dal 2015 & Amazon Echo/Alexa - "effetto Alexa" ampiamente documentato & alta \\
Italia & Erica & 2001-02 & 863(1999)->416(2002), quasi dimezzato & Delitto di Novi Ligure (Erika De Nardo uccide madre e fratello) & alta \\
\bottomrule
\end{tabular}
\end{table*}


\subsection{Un controllo di realtà: gli stati sono tutti d'accordo?}

Per i casi-studio più solidi (sei positivi, tre negativi, tutti americani poiché non esistono dati
regionali italiani comparabili), abbiamo verificato se il salto o il crollo osservato a livello
nazionale fosse effettivamente diffuso su tutto il paese, o concentrato in poche zone
(Figura~\ref{fig:state_concentration}, Tabella~\ref{tab:state_concentration}). La logica è semplice:
un vero evento mediatico nazionale dovrebbe muovere la maggior parte degli stati nella stessa
direzione; un effetto concentrato in uno o due stati soli sarebbe un campanello d'allarme per una
causa regionale o, peggio, per una coincidenza statistica. Sette casi su nove superano il 90\% di
stati concordanti (Jaime al 100\%, Kobe e Alexa al 97\%), il che ci dà una certa fiducia che questi non
siano artefatti isolati. Le due eccezioni sono interessanti di per sé: Jaslene (79\%) è compatibile
con un'adozione inizialmente più forte nelle aree a maggiore presenza ispanica, coerente con il fatto
che fu la prima vincitrice ispanica del programma; Hillary (74\%) è compatibile con una polarizzazione
politica che, ragionevolmente, non ha colpito ogni stato allo stesso modo.

\begin{figure}[t]
\centering
\includegraphics[width=\linewidth]{figures/fig8_state_concentration.png}
\caption{Quota di stati concordanti con l'effetto nazionale, per caso-studio.}
\label{fig:state_concentration}
\end{figure}

\begin{table*}[t]
\centering
\scriptsize
\caption{Uniformità geografica dei casi-studio RQ2 (dati per stato, solo USA)}
\label{tab:state_concentration}
\begin{tabular}{l c c c c r c r c r c r}
\toprule
\textbf{nome} & \textbf{sesso} & \textbf{anno\_evento} & \textbf{tipo} & \textbf{stati\_con\_dati} & \textbf{stati\_concordanti} & \textbf{percento\_stati\_concordanti} & \textbf{rapporto\_mediano\_stati} & \textbf{stato\_top} & \textbf{rapporto\_stato\_top} & \textbf{stato\_minimo} & \textbf{rapporto\_stato\_minimo} \\
\midrule
Tammy & F & 1958 & positivo & 51 & 47 & 92.2 & 2.741 & NC & 59.667 & DC & 0.933 \\
Jaime & F & 1976 & positivo & 45 & 45 & 100.0 & 8.053 & PA & 44.083 & DC & 3.667 \\
Devante & M & 1992 & positivo & 11 & 11 & 100.0 & 12.333 & FL & 21.800 & CA & 3.864 \\
Mariah & F & 1991 & positivo & 46 & 45 & 97.8 & 6.167 & NY & 23.083 & ND & 1.167 \\
Nevaeh & F & 2001 & positivo & 7 & 7 & 100.0 & 7.333 & CA & 11.091 & IL & 4.333 \\
Jaslene & F & 2008 & positivo & 19 & 15 & 78.9 & 1.667 & PA & 3.667 & CT & 1.000 \\
Hillary & F & 1993 & negativo & 43 & 32 & 74.4 & 0.571 & LA & 1.031 & MN & 0.292 \\
Kobe & M & 2004 & negativo & 37 & 36 & 97.3 & 0.471 & MO & 0.815 & KY & 0.222 \\
Alexa & F & 2020 & negativo & 37 & 36 & 97.3 & 0.414 & ID & 1.000 & OH & 0.250 \\
\bottomrule
\end{tabular}
\end{table*}


\subsection{Le storie in comune tra i due paesi}

Infine, quattro casi che permettono un confronto diretto tra USA e Italia sullo stesso evento.
\textbf{Celine Dion} è forse il più istruttivo: la sua diagnosi di sindrome della persona rigida
(annunciata a dicembre 2022) e il documentario che ne è seguito (giugno 2024) hanno prodotto una
crescita graduale del nome negli USA (da 614 a 1.466 nascite tra il 2018 e il 2025) ma un'impennata
molto più marcata in Italia (fino a 6 volte il livello di base): stessa causa, magnitudine molto
diversa, probabilmente perché il nome partiva da una base molto più piccola in Italia.
\textbf{Elsa}, dopo \emph{Frozen}, è reale in entrambi i paesi ma arriva in Italia con circa un anno di
ritardo (il film uscì lì a dicembre 2013) e con un'ampiezza molto più contenuta. \textbf{Khaleesi}
resta invece un fenomeno puramente americano: non compare mai nei dati italiani catturati, nemmeno
negli anni a copertura completa. \textbf{Isabella}, dopo \emph{Twilight}, è il caso più ambiguo: cresce
in entrambi i paesi, ma in Italia la crescita comincia \emph{prima} dell'uscita della saga e prosegue
senza scosse anche dopo, probabilmente perché Isabella è già un nome classico italiano, e quindi le
due curve, per quanto visivamente simili, raccontano storie causali diverse.

\subsection{Una tabella esaustiva, non solo i casi scelti a mano}

Per non dare l'impressione di aver scelto solo gli esempi più belli da raccontare, abbiamo anche
generato una tabella puramente meccanica di ogni nome/anno che supera una soglia di crescita fissata a
priori (USA: rapporto ${\geq}6\times$ e conteggio finale ${\geq}2.000$; Italia: rapporto
${\geq}2{,}5\times$ e conteggio ${\geq}150$), per un totale di 87 righe (48 USA, 39 Italia), senza
alcuna selezione manuale. Le tredici storie curate qui sopra sono un sottoinsieme di questa base più ampia,
non un'eccezione ad essa.

\section{Divergenza nei Ruoli di Genere Anagrafici: Nomi Gender-Neutral (RQ3)}
\label{sec:rq3}

La terza domanda di ricerca (RQ3) analizza se il processo di individualizzazione e l'evoluzione delle norme di genere abbiano portato ad un aumento dei \textbf{nomi gender-neutral (unissex)}, ovvero attribuiti sia a nati maschi che femmine ($0.10 \le p_F \le 0.90$).

Come sintetizzato nella Tabella~\ref{tab:gender_neutral} e nella Figura~\ref{fig:gender_neutral}, si osserva una drastica divergenza strutturale tra i due paesi:

\begin{enumerate}
    \item \textbf{Stati Uniti (USA)}: La quota di nascite con nomi neutri rispetto al sesso è cresciuta dal \textbf{$7{,}31\%$} nel 1999 all'\textbf{$8{,}81\%$} nel 2024 (test di Mann-Kendall: $\tau = 0.526, p = 1.79 \times 10^{-4}$, altamente significativo). Nomi come \emph{Avery}, \emph{Riley}, \emph{Jordan}, \emph{Charlie}, \emph{Taylor}, \emph{Peyton}, \emph{Morgan} e \emph{Quinn} sono entrati stabilmente nell'uso comune.
    \item \textbf{Italia (IT)}: La quota di nascite con nomi unissex rimane estremamente limitata, passando da uno \textbf{$0{,}06\%$} nel 1999 allo \textbf{$0{,}22\%$} nel 2024 (test di Mann-Kendall: $\tau = 0.031, p = 0.842$, non significativo).
\end{enumerate}

\begin{table*}[t]
\centering
\footnotesize
\caption{Nomi Gender-Neutral / Unissex, USA vs Italia, 1999-2024 (RQ3)}
\label{tab:gender_neutral}
\begin{tabular}{>{\raggedright\arraybackslash}p{2.3cm} c c c c c >{\raggedright\arraybackslash}p{2.2cm}}
\toprule
\textbf{Paese} & \textbf{Quota 1999 (\%)} & \textbf{Quota 2024 (\%)} & \textbf{Trend} & \textbf{Tau} & \textbf{p-value} & \textbf{Esito} \\
\midrule
Stati Uniti (USA) & 7.31\% & 8.81\% & Crescente & 0.526 & 1.7892e-04 & Significativo \\
Italia (IT) & 0.06\% & 0.22\% & Stabile & 0.031 & 8.4239e-01 & Non significativo \\
\bottomrule
\end{tabular}
\end{table*}


\begin{figure*}[t]
\centering
\includegraphics[width=0.9\linewidth]{figures/fig9_gender_neutral.png}
\caption{Evoluzione dei nomi Gender-Neutral (Unissex) 1999-2024: quota di nascite totali (sinistra) e conteggio nomi unissex distinti (destra).}
\label{fig:gender_neutral}
\end{figure*}

Questa marcata differenza riflette sia vincoli di natura normativa (l'art. 35 del DPR 396/2000 in Italia ha storicamente imposto la corrispondenza del nome al sesso biologico dell'atto di nascita) sia una maggiore persistenza delle convenzioni di genere tradizionali nella scelta anagrafica italiana.

\section{Discussione Finale}
\label{sec:discussione}

Il quadro che emerge da questo lavoro racconta una storia coerente attraverso le tre domande di ricerca (diversità di lungo periodo RQ1, impatti mediatici RQ2, ed evoluzione delle norme di genere RQ3).
In tutti e tre i casi, il filo conduttore è lo stesso: la scelta del nome di un figlio, per quanto intima, resta sistematicamente esposta a forze culturali che vanno ben oltre il singolo nucleo familiare, tra cui grandi cambiamenti generazionali di valori (la diversificazione di lungo periodo, coerente con la tesi di Twenge et al. sull'aumento dell'individualismo), la globalizzazione dei media (la convergenza USA-Italia, con nomi anglo-americani come Liam e Noah che attraversano l'Atlantico), l'impatto causale della finzione narrativa e della musica pop, ed il superamento progressivo dei ruoli tradizionali di genere negli USA rispetto all'Italia (RQ3).

Un punto che vale la pena riprendere esplicitamente qui è quello, già accennato in Metodologie, sull'asimmetria tra le due infrastrutture statistiche nazionali. Gli Stati Uniti pubblicano l'intera distribuzione dei nomi in un archivio aperto, aggiornato annualmente, ricco dal 1880. L'Italia non pubblica nulla del genere: i dati usati in questo lavoro esistono solo perché raggiungibili tramite un servizio web non documentato, pensato per un tool di consultazione singola, non per la ricerca. Non è un dettaglio tecnico neutro, ma a suo modo un dato culturale: riflette due approcci molto diversi alla trasparenza dei dati pubblici, un tema che meriterebbe una riflessione a sé.

\section{Conclusioni}

Questo lavoro ha risposto in modo quantitativo e rigoroso al complesso di quesiti sull'evoluzione dell'onomastica transatlantica. La diversità cresce in entrambi i paesi; l'Italia resta più concentrata, ma il divario femminile si sta chiudendo; le scelte dei due paesi si stanno avvicinando in modo tracciabile; gli eventi cinematografici e musicali producono shock reali con adozione strutturale permanente a 5 anni; ed infine i ruoli di genere evidenziano un aumento di nomi unissex negli USA che stacca nettamente la rigidità anagrafica italiana.

\section*{Repository e riproducibilità}

Codice, dati grezzi ed elaborati, figure e tabelle sono disponibili nel repository del progetto:
\url{https://github.com/matteraggi/HDS_evolution_of_naming}.

% Fonti verificate via ricerca web il 2026-08-17 (non solo a memoria). Correzione rispetto
% a una versione precedente di questa bozza: Mann (1945) riportava le pagine 163--171 a
% memoria, sbagliato - il numero corretto (confermato da Cambridge Core, Econometric
% Society, Semantic Scholar) è 245--259, corretto qui sotto.
\begin{thebibliography}{00}

\bibitem{mann1945} H.~B. Mann, ``Nonparametric Tests Against Trend,'' \emph{Econometrica}, vol.~13,
no.~3, pp.~245--259, 1945. doi: 10.2307/1907187.

\bibitem{kendall1975} M.~G. Kendall, \emph{Rank Correlation Methods}, 4th ed. London: Charles
Griffin, 1975.

\bibitem{wilcoxon1945} F.~Wilcoxon, ``Individual Comparisons by Ranking Methods,'' \emph{Biometrics
Bulletin}, vol.~1, no.~6, pp.~80--83, 1945. doi: 10.2307/3001968.

\bibitem{lieberson2000} S.~Lieberson, \emph{A Matter of Taste: How Names, Fashions, and Culture
Change}. New Haven: Yale University Press, 2000. ISBN 0-300-08385-8.

\bibitem{lieberson1992} S.~Lieberson and E.~O. Bell, ``Children's First Names: An Empirical Study of
Social Taste,'' \emph{American Journal of Sociology}, vol.~98, no.~3, pp.~511--554, 1992.
doi: 10.1086/230048.

\bibitem{twenge2010} J.~M. Twenge, E.~M. Abebe, and W.~K. Campbell, ``Fitting In or Standing Out:
Trends in American Parents' Choices for Children's Names, 1880-2007,'' \emph{Social Psychological and
Personality Science}, vol.~1, no.~1, pp.~19--25, 2010. doi: 10.1177/1948550609349515.

\bibitem{twenge2016} J.~M. Twenge, L.~Dawson, and W.~K. Campbell, ``Still Standing Out: Children's
Names in the United States During the Great Recession and Correlations with Economic Indicators,''
\emph{Journal of Applied Social Psychology}, vol.~46, pp.~663--670, 2016. doi: 10.1111/jasp.12409.

\bibitem{ogihara2025a} Y.~Ogihara, ``Popularity and Diversity: The Negative Relationship in Baby
Names in the United Kingdom,'' \emph{F1000Research}, vol.~14, p.~424, 2025.
doi: 10.12688/f1000research.162476.1.

\bibitem{ogihara2025b} Y.~Ogihara, ``Name Uniqueness and the Rise of Individualism in the Western
Hemisphere (1500-2000),'' \emph{Current Research in Ecological and Social Psychology}, vol.~9, 2025.

\bibitem{ons2015} UK Office for National Statistics, T.~Davy and R.~Lewis, ``10 Pop Culture
Influences on Baby Names: Game of Thrones, Marvel, Frozen and More,'' 17 August 2015. [Online].
Available: \url{https://www.ons.gov.uk/peoplepopulationandcommunity/birthsdeathsandmarriages/livebirths/articles/10popcultureinfluencesonbabynamesgameofthronesmarvelfrozenandmore/2015-08-17}

\end{thebibliography}

\end{document}
